\documentclass[11pt,a4paper]{article}
\usepackage[T1]{fontenc}
\usepackage[utf8]{inputenc}
\usepackage{lmodern}
\usepackage[margin=28mm]{geometry}
\usepackage{amsmath,amssymb,amsthm,mathtools}
\usepackage{booktabs,array,microtype}
\usepackage{url}
\usepackage[colorlinks=true,linkcolor=blue,citecolor=blue,urlcolor=blue]{hyperref}
\newtheorem{theorem}{Theorem}
\newtheorem{lemma}[theorem]{Lemma}
\newtheorem{proposition}[theorem]{Proposition}
\theoremstyle{remark}
\newtheorem{remark}[theorem]{Remark}
\newcommand{\eps}{\varepsilon}
\newcommand{\PP}{\mathrm{P}}
\newcommand{\QQ}{\mathrm{Q}}
\newcommand{\Lang}[1]{\mathcal{#1}}
\newcommand{\Z}{\mathbb{Z}}
\DeclareMathOperator{\Av}{Av}
\title{A direct first-return enumeration of\\$(2,2)$-stack-sortable permutations}\hypersetup{pdfauthor={Xiaojian Zeng},pdftitle={A direct first-return enumeration of (2,2)-stack-sortable permutations}}
\author{Xiaojian Zeng\\[2pt]
  \small PBC School of Finance, Tsinghua University, 100190\\
  \small \href{mailto:xiaojianzeng@foxmail.com}{xiaojianzeng@foxmail.com}}
\date{September 12, 2026}

\begin{document}
\maketitle
\begin{abstract}
A $(2,2)$-stack permits insertion and removal at either of its first
two available positions. We derive the generating function of its sortable
input permutations directly from these operations, answering Problem~5.11
of Vatter. An eager deterministic sorter chooses one operation word for each
sortable input. The comparisons imposed by its pushes form an ordered
forest, and a leaf/nonleaf bit on each live token preserves every remaining
comparison obligation. The resulting bit-stack language has an unambiguous
first-return decomposition with two endpoint states. Its equations give
$F=1+xAF$ and $A=1+xA(4A-2F-1)$, and hence the previously known cubic
$2xF^3-(2x+3)F^2-(x-7)F-4=0$, with $F(0)=1$.
The contribution is the operation-level bijection and decomposition;
the polynomial and counting sequence were already known. Exact finite
checks accompany the proof, which applies to every input length.
\end{abstract}

\section{Model and contribution}

An input permutation of $[n]=\{1,\ldots,n\}$ is read from left to right.
The stack and output are initially empty. Reading the stack from top to
bottom, a push inserts the next unread entry at position~1 or~2, when that
position is available, and a pop removes the entry at position~1 or~2,
when present, and appends it to the output. In particular, an empty stack
has only the first insertion position. An input is \emph{sortable} if some
legal operation sequence consumes it and produces $1,\ldots,n$, leaving
the stack empty. Let $u_n$ count sortable inputs once each, regardless of
how many sequences sort them, and set
\[
 F(x)=\sum_{n\geq0}u_nx^n\in\Z[[x]],\qquad u_0=1.
\]
We write $\PP_j$ and $\QQ_j$ for a push and a pop at rank~$j$.

Atkinson studied the more general $(r,s)$-stacks, with insertion in the
first $r$ positions and removal from the first $s$ positions, including
forbidden-subsequence characterizations for $r=s=2$~\cite{Atkinson1998}.
Vatter~\cite[Theorem~5.9 and the paragraph before Problem~5.11]{Vatter2026}
records the eight-element basis and attributes a derivation of the cubic
below to Pantone, using Combinatorial Exploration from that basis.
The framework of Albert et al.~\cite{AlbertEtAl2026} derives counting
equations from combinatorial specifications. Vatter's Problem~5.11 asks
instead for a derivation from the stack's structure. We provide such a
derivation through the comparisons required by a canonical sorting run.

\begin{theorem}\label{thm:main}
For every $n\geq0$, sortable input permutations of $[n]$ are in bijection
with complete words having $n$ pushes in the following bit-stack machine,
subject to exclusion of the adjacent factor $\PP_2\QQ_1$:
\begin{equation}\label{eq:bits}
\begin{array}{rcll}
 \PP_1:&\eps\longmapsto0,&\quad aT\longmapsto1aT,\\
 \PP_2:&aT\longmapsto10T,\\
 \QQ_1:&aT\longmapsto T,\\
 \QQ_2:&0bT\longmapsto0T.
\end{array}
\end{equation}
Here $a,b\in\{0,1\}$, $T$ is any bit suffix, and complete words start and
end with the empty stack. Labeling tokens by increasing pop order and
reading their labels in push order gives the corresponding input.
The first-return specification in Section~\ref{sec:grammar} is complete
and unambiguous. Its generating functions satisfy
\begin{equation}\label{eq:reduced}
 F=1+xAF,\qquad A=1+xA(4A-2F-1),
\end{equation}
where $A$ counts the positive-floor paths defined there. Consequently,
$F$ is the unique formal power series with constant term~1 satisfying
\begin{equation}\label{eq:cubic}
 2xF^3-(2x+3)F^2-(x-7)F-4=0.
\end{equation}
\end{theorem}

Equation~\eqref{eq:cubic} and its coefficients are known; they also appear
in the source-author OEIS entry A393395~\cite{OEIS393395}. No pattern basis
or preassigned polynomial enters the proof below. We work throughout
with sorting an arbitrary input into increasing order, so no change from
sorting to generation, or associated permutation symmetry, is needed.

\section{A canonical sorting word}

At any state of a successful run, let $k$ be the next required output.
We first establish two exchange observations for the original machine.

\begin{lemma}\label{lem:eager}
If a state admits a successful continuation and $k$ is in one of the
first two stack positions, popping $k$ immediately preserves this property.
\end{lemma}
\begin{proof}
Before its pop of $k$, a successful continuation consists entirely of
pushes. Pop $k$ now and replay those pushes with $k$ absent. For an original
insertion slot $j$, use $j-1$ if $k$ lay strictly before that slot, and
use $j$ otherwise. The new slot is available and is at most~2. The only
possible reduction is from slot~2 to slot~1, so slot~0 never occurs.
After omitting the original pop of $k$, the two stacks, unread inputs,
and outputs coincide. The remaining continuation is unchanged.
\end{proof}

\begin{lemma}\label{lem:pair}
Fix an unread input and stack suffix $T$. For $a<b$, a successful
continuation from stack $(b,a,T)$ implies one from $(a,b,T)$, with the
same output already produced.
\end{lemma}
\begin{proof}
Follow a successful continuation from $(b,a,T)$ until it removes $a$,
interchanging $a$ and $b$ in the second run. Neither token is removed
earlier, because output is increasing. Every operation on another token
has the same rank in both stacks. The relative order of surviving tokens
is unchanged by rank operations, so $b$ still precedes $a$ in the first
stack. Since $a$ is accessible when removed, that stack must then begin
with $(b,a)$. Remove $a$ at rank~1 in the second run, instead of rank~2
in the first. The remaining stacks agree, and the continuation can be
copied without further change.
\end{proof}

Define the deterministic sorter $G$ as follows. Pop $k$ whenever it is
in the first two positions. Otherwise push the next unread input $a$,
using rank~1 if the stack is empty or $a$ is smaller than the old top $b$,
and rank~2 if $b<a$. If no such pop or push is possible, accept only if
both the input and stack are exhausted, and reject otherwise.

\begin{proposition}\label{prop:canonical}
The sorter $G$ accepts exactly the sortable inputs. Every accepted input
has exactly one $G$ word, and complete $G$ words determine their inputs
uniquely.
\end{proposition}
\begin{proof}
Soundness follows from the legal operations and the increasing output.
For completeness, induct on the number of operations remaining in a
successful continuation. If an eager pop is possible, Lemma~\ref{lem:eager}
preserves success. Otherwise a push is forced. On a nonempty stack its
two possible ranks create the two orders of the same top pair, and
Lemma~\ref{lem:pair} shows that the smaller-top order chosen by $G$
preserves success. On an empty stack the push is unique. This proves
completeness, including the terminal state.

More generally, take any complete legal rank word, assigning a distinct
token to each push. Label the token removed at the $j$th pop by~$j$.
Reading these labels in push order gives the unique permutation for which
the word produces increasing output. Restricting this reconstruction to
$G$ words is injective because $G$ is deterministic, and it is surjective
onto the sortable inputs by completeness. Each push accounts for exactly
one input entry.
\end{proof}

\section{Comparisons, forests, and bits}\label{sec:forest}

Use the reconstructed labels of a complete rank word. A nonempty-stack
$\PP_1$ is the greedy choice precisely when its new token is output
before the old top; $\PP_2$ is the greedy choice precisely when the old
top is output before the new token. Record this obligation as a directed
edge from the upper token of the new top pair to the lower token.
An empty-stack push creates no edge. Call a word \emph{comparison-valid}
if the source of every edge is popped before its target.

\begin{lemma}[Peak criterion]\label{lem:peak}
A comparison-valid complete word is a $G$ word if and only if each
adjacent push-pop pair is $\PP_1\QQ_1$.
\end{lemma}
\begin{proof}
In a $G$ run, consider a push immediately followed by a pop. The required
token was inaccessible before the push. Insertion cannot move an old
token closer to the top, so the pop removes the new token. This token
is the next required output and is smaller than every resident token;
it was therefore pushed at rank~1 and is popped there.

Conversely, consider a moment immediately before any push in a complete
word satisfying the criterion. The next pop follows a nonempty run of
pushes. Its last push and first pop form $\PP_1\QQ_1$, so the token
removed was not yet present at the chosen moment. Under reconstruction
it is precisely the next required output. Thus an eager pop was not
available. Comparison-validity supplies $G$'s push choice. At every pop,
reconstruction makes the removed token the next required output.
Every operation therefore agrees with $G$.
\end{proof}

The comparison graph can be maintained using only the unpopped tokens.
When a token is removed, delete it and all its incident edges. An ordered
rooted forest is read in preorder, visiting a root and then its ordered
child subtrees, before moving to the next root.

\begin{lemma}\label{lem:forest}
After any prefix whose pops respect the comparison edges, the live
graph is an ordered rooted forest whose preorder is the stack order.
The first stack token is always a root. The second token, when present,
is a root exactly when the first token is a leaf.
\end{lemma}
\begin{proof}
The initial forest is empty. A push into an empty stack creates a
singleton tree. A nonempty $\PP_1$ creates a root whose only child is
the previous first root; later trees are unchanged. A $\PP_2$ adds a new
first child to the first root. Both updates preserve the claimed preorder.

The first token is a root, so $\QQ_1$ removes it and promotes its children
to roots in their existing order. If that root is a leaf, the second
token is the next root, and $\QQ_2$ removes the latter and promotes its
children after the surviving first root. If the first root is not a leaf,
the second token is its first child. Removing that child would violate
the edge from the first root. These are all available pop ranks, proving
the invariant by induction.
\end{proof}

This argument gives necessity and sufficiency: a comparison-valid word
cannot make the forbidden second-rank pop, and removing only roots
respects every comparison ever created. Store bit~0 on a leaf and bit~1
on a nonleaf. A surviving vertex cannot lose a child, since a child has
an unmet incoming edge until its parent is removed. Thus one never needs
to count the children or change a surviving token's bit from~1 to~0.
The forest updates are exactly~\eqref{eq:bits}. In particular, $\PP_2$
marks the old top as a nonleaf and inserts a new leaf in second position.

A push into a nonempty stack leaves top bit~1, which prevents a following
$\QQ_2$. A push into an empty stack leaves only one token, so again
$\QQ_2$ cannot follow it. The bit rules therefore prohibit
$\PP_1\QQ_2$ and $\PP_2\QQ_2$ automatically. Excluding the single factor
$\PP_2\QQ_1$ leaves exactly the peaks in Lemma~\ref{lem:peak}.
Together with Proposition~\ref{prop:canonical}, this proves the bijection
asserted in Theorem~\ref{thm:main}.

\section{The first-return specification}\label{sec:grammar}

We now enumerate rank words of the abstract bit machine, always excluding
$\PP_2\QQ_1$. Height means the number of stack tokens. A
\emph{positive-floor path} starts at positive height with top bit $a$ and
suffix $T$, ends at the same height, and never goes below that height.
A \emph{first-return excursion} is nonempty and stays strictly above
its initial height except at its endpoints.

The suffix $T$ is protected throughout such paths. During a first-return
excursion, before its closing step there are at least two tokens above
$T$. No permitted rank can reach $T$ without ending the excursion first.
After the opening push, the second token and $T$ are similarly protected
throughout the interior, whose floor is one unit higher. Pushes inspect
or modify only the active top bit. It follows inductively over successive
excursions that the language and final top bit do not depend on the
contents of $T$. The initial top token may nevertheless be replaced.
The interface must therefore record its bit, not its identity.

Opening $\PP_1$ from top bit $a$ creates active top bit~1 and stores $a$
second. Opening $\PP_2$ creates active top bit~1 and stores~0 second.
The interior is a positive-floor path starting at bit~1. Closing $\QQ_1$
restores the stored bit. Closing $\QQ_2$ requires the interior to end at
bit~0; it deletes the stored token and retains that~0. An empty interior
with opening $\PP_2$ and closing $\QQ_1$ is excluded.

\begin{lemma}\label{lem:endpoint}
A positive-floor path starting at bit~0 can end only at bit~0. With
the endpoint languages
\[
 \Lang A:0\longrightarrow0,\qquad
 \Lang B:1\longrightarrow0,\qquad
 \Lang C:1\longrightarrow1,
\]
where $\Lang A,\Lang C$ contain $\eps$ and $\Lang B$ does not, define
\begin{align}
 \Lang E&=\PP_1(\Lang B\sqcup\Lang C)\QQ_1,\label{eq:E}\\
 \Lang D&=\PP_1\Lang B\QQ_2
 \sqcup\PP_2\Lang B\QQ_1
 \sqcup\PP_2(\Lang C\setminus\{\eps\})\QQ_1
 \sqcup\PP_2\Lang B\QQ_2.\label{eq:D}
\end{align}
Then all positive-floor paths have the unambiguous specification
\begin{align}
 \Lang A&=\{\eps\}\sqcup\Lang E\Lang A\sqcup\Lang D\Lang A,
       \label{eq:A}\\
 \Lang C&=\{\eps\}\sqcup\Lang E\Lang C,\label{eq:C}\\
 \Lang B&=\Lang E\Lang B\sqcup\Lang D\Lang A.\label{eq:B}
\end{align}
If $\Lang L$ is the complete empty-stack language, then
\begin{equation}\label{eq:L}
 \Lang L=\{\eps\}\sqcup\PP_1\Lang A\QQ_1\Lang L.
\end{equation}
Here $\sqcup$ denotes disjoint union and products mean concatenation.
\end{lemma}
\begin{proof}
The assertion about an initial~0 follows by induction on path length
from the opening and closing cases just described. For one excursion,
a closing $\QQ_1$ restores either the initial~0 or the~0 stored by
$\PP_2$, and a closing $\QQ_2$ retains~0. Apply the induction hypothesis
to the shorter remainder after this first excursion.

For initial bit~1, excursions of type $\Lang E$ preserve~1, while the
four alternatives of $\Lang D$ return~0. For initial bit~0 both types
return~0. More explicitly, the exhaustive split is
\begin{center}
\begin{tabular}{lll}
\toprule
Opening and closing & Interior & Final bit\\
\midrule
$\PP_1\ \cdots\ \QQ_1$ & $\Lang B\sqcup\Lang C$ & initial bit\\
$\PP_1\ \cdots\ \QQ_2$ & $\Lang B$ & $0$\\
$\PP_2\ \cdots\ \QQ_1$ & $\Lang B\sqcup(\Lang C\setminus\{\eps\})$ & $0$\\
$\PP_2\ \cdots\ \QQ_2$ & $\Lang B$ & $0$\\
\bottomrule
\end{tabular}
\end{center}
Every nonempty path splits into its first-return excursion and remaining
path. This gives~\eqref{eq:A}--\eqref{eq:B}. A nonempty empty-stack path
starts with $\PP_1$, creating bit~0. Its first component ends with
$\QQ_1$, the only pop available at height~1, and has interior in
$\Lang A$. This proves~\eqref{eq:L}.

For unambiguity, the height profile determines the first-return time,
the interior, and the remaining suffix uniquely. The first and last
operation ranks distinguish every excursion alternative except the two
$\PP_2\cdots\QQ_1$ terms in~\eqref{eq:D}. Their interior endpoint,
which is uniquely determined by the bit transitions, distinguishes
$\Lang B$ from $\Lang C$. Induction gives a unique parse.
Conversely, every production obeys the bit transitions, floor restriction,
and forbidden-factor condition. Joining nonempty excursions creates only
a pop-push boundary. The only problematic empty interior was expressly
removed in~\eqref{eq:D}. Thus no extra words or multiple parses arise.
\end{proof}

\begin{remark}
Auxiliary initial bit~1 is a state of the abstract bit machine; it need
not be realized by a physical forest consisting of a single nonleaf.
Its suffix can be supplied without changing its path language. The
complete language starts at the actual empty state, where the forest
correspondence already proves realizability. The bit stack has unbounded
height; the finite interface above is a boundary description, not a
claim that the entire sorter has finitely many states.
\end{remark}

\section{Formal power series}\label{sec:series}

Weight each push by $x$, and give pops weight~1. Use $A,B,C$ for the
ordinary generating functions of $\Lang A,\Lang B,\Lang C$, and let
$e,d$ count the first-return languages $\Lang E,\Lang D$. The translation
of disjoint unions and unambiguous products into sums and products is
the symbolic method~\cite[Chapter~I]{FlajoletSedgewick2009}; the role of
productivity in recursive specifications is treated explicitly by
Albert et al.~\cite[Sections~3.3, 4 and~5.2]{AlbertEtAl2026}.
Here the translation can be read directly from the displayed grammar:
\begin{align}
 e&=x(B+C),& d&=x(3B+C-1),\label{eq:ed}\\
 A&=1+(e+d)A,& C&=1+eC,\label{eq:AC}\\
 B&=eB+dA,& F&=1+xAF.\label{eq:BF}
\end{align}
The subtraction of~1 in $d$ removes precisely the excluded empty
$\Lang C$ interior. Every excursion contributes at least one push;
thus, with $A(0)=C(0)=F(0)=1$ and $B(0)=0$, coefficient $n$ in these
equations depends only on smaller-size coefficients. This supplies
productivity and finite coefficients directly in this instance.

Since $e(0)=d(0)=0$, all the following inverses exist in $\Z[[x]]$:
\[
 A=\frac1{1-e-d},\qquad C=\frac1{1-e},\qquad
 B=dAC=A-C.
\]
Hence $B+C=A$, so $e=xA$ and
\[
 C=\frac1{1-xA}=F.
\]
Substituting $B=A-F$ into~\eqref{eq:ed}--\eqref{eq:BF} gives
exactly~\eqref{eq:reduced}.

For an elimination without division by $x$, set
\[
 R_1=F-1-xAF,\qquad R_2=A-1-xA(4A-2F-1).
\]
Direct expansion yields
\begin{align*}
 &2xF^3-(2x+3)F^2-(x-7)F-4\\
 &\qquad=xF^2R_2+
 \bigl(-4xAF+2xF^2+xF-3F+4\bigr)R_1.
\end{align*}
Both residuals vanish, proving~\eqref{eq:cubic}. To check the branch
uniquely, write $F=1+H$ with $H(0)=0$. The left side becomes
\[
 H-3H^2+x(-1+H+4H^2+2H^3).
\]
At each positive degree the next coefficient of $H$ occurs with
multiplier~1, and every other term depends only on earlier coefficients.
Thus there is at most one formal solution with $F(0)=1$, and the grammar
constructs it. This completes the proof of Theorem~\ref{thm:main}.

\section{Exact checks and scope}\label{sec:checks}

The accompanying programs use exact integers and finite word sets.
To check the compression against the original machine, a separate
implementation reconstructs the unrestricted legal sorting search without
using either exchange lemma. Separately rebuilding the finite structure
in an exact verifier is also used by Zeng et al.~\cite{ZengEtAl2026}
for ternary codes (Section~4, Proposition~1): their two implementations
construct the compatibility graph independently. In the present problem,
the unrestricted machine and the bit language play those distinct
verification roles.

The unrestricted search and a separately implemented $G$ agree on all
$46\,234$ permutations across $0\leq n\leq8$. In that range, complete
bit words agree objectwise with canonical words, reconstruction gives
the same input, and every complete word has exactly one grammar parse.
The initial counts are
\begin{center}
\small
\begin{tabular}{c*{10}{r}}
\toprule
$n$ & 0&1&2&3&4&5&6&7&8&9\\
\midrule
$u_n$ &1&1&2&6&24&116&628&3636&21956&136428\\
\bottomrule
\end{tabular}
\end{center}
These agree with A393395~\cite{OEIS393395}. The $n=9$ check compares $G$
on all $362\,880$ inputs with all $136\,428$ complete bit words, including
reconstructed inputs and multiplicities. It does not extend the
separately implemented unrestricted sorting search beyond $n=8$.

The recursive word generator also checks all auxiliary endpoint languages
through eight pushes. A separate parser and bit implementation check the
positive-floor languages through seven pushes. Through five pushes,
$13\,860$ replays test every protected bit suffix of lengths zero through
three, including preservation of its token identities. An additional
attack includes rejected words: it enumerates all legal rank words at
floor~0 through six pushes and at floor~1 through five pushes, totaling
$107\,848$ words. The grammar's parse multiplicity always equals the
appropriate bit-language membership indicator. No discrepancy occurs in
these stated ranges. Independent coefficient recursions agree through
degree~40, and exact polynomial expansion checks the residual identity.

These checks detect finite implementation errors; the forest invariant,
endpoint argument, and unambiguous grammar establish the result at all
lengths. The result concerns only the $(2,2)$-stack. In particular, it
makes no assertion about general $(r,s)$-stacks or the novelty of the
elementary greedy normalization by itself. The inspected primary
comparison identifies the basis-derived enumeration as the closest
documented result. A bounded literature search completed on September~12,
2026 found no equivalent direct derivation; it does not establish global
priority. The historical account of Pantone's enumeration is Vatter's
report, rather than an inspection of Pantone's computation.

\paragraph{Computational and AI assistance.}
OpenAI Codex assisted with proof development, literature searches,
implementation of the exact checks, a separate model-based mathematical
review, and manuscript preparation. The review was a separate model
context, not external peer review. Language-model output and finite
agreement are not substitutes for the proofs above. The symbolic check
used SymPy; the combinatorial checks used Python integer arithmetic.
No proof assistant was used. The supplementary documentation specifies
the programs, finite ranges, and reproducible commands.

\bibliographystyle{plainurl}
\bibliography{references}
\end{document}
